\section{科研心态与基本规范}

\subsection{科研与上课的差别}
做科研通常没有现成教科书可逐步照做，主要依靠文献检索、批判性阅读、实验验证与反复修改。进展慢、卡住、被拒稿都很常见；关键是建立可持续的习惯，而不是追求短期「看起来很忙」。

\subsection{能力与态度}
结合组内学长与学术交流共识，入门阶段尤其需要：
\begin{itemize}
  \item \textbf{主动性}：主动查文献、主动问清任务边界、事事有回音；
  \item \textbf{自主学习}：会用检索与社区资源，并能归纳总结；
  \item \textbf{抗挫折}：论文读不懂、代码跑不通、实验不涨点是常态；
  \item \textbf{时间管理}：期刊周期常要数月，需提前为写作与返修留时间；
  \item \textbf{身心健康}：科研只是生活的一部分；适当运动有助于长期效率。
\end{itemize}

\subsection{常见误区}
\begin{enumerate}
  \item \textbf{「顶会/顶刊文章一定都对」}：带批判性阅读；A 类未必处处优于 B 类，没有工作是完美的。
  \item \textbf{「调参总能调出好结果」}：好模型不应只靠黑箱撞运气；单组参数极好而其他组都崩，往往意味着错误或偶然性，需警惕。
  \item \textbf{「模型越复杂越好」}：优先抓住核心创新与可解释动机，避免为复杂而复杂。
  \item \textbf{「不清楚自己解决了什么问题」}：投稿前反复自问：真问题是什么？技术贡献是什么？读者能带走什么？切忌编造问题。
\end{enumerate}

\subsection{选题与问题价值（交流纪要精华）}
确定方向前，建议先想清楚：
\begin{itemize}
  \item 问题为什么重要？现有方法为什么解决不好？
  \item 自己为什么有机会解决（数据、算力、团队积累）？
  \item 能否逐步形成一条研究主线，而不是互不相关的零散文章？
\end{itemize}
有价值的问题往往来自长期看数据、做实验、接触真实场景，而不是简单从别人 limitation 里「抄一个缺口」。换方法、换技术点可以，但尽量避免频繁更换大方向。

\subsection{何时收尾、如何面对拒稿}
\begin{itemize}
  \item 论文没有绝对完美版。若问题讲清、方法合理、关键实验稳定、结论可支撑，就应整理投稿，而不是无限打磨。
  \item 拒稿是常态。有道理的意见认真改；不合理的意见不必过度内耗。rebuttal 要判断工作本身是否够强、负面意见是否可正面回应。
\end{itemize}

\subsection{硕士阶段可落地的启发}
\begin{enumerate}
  \item 困难是常态，不要因短期无成果否定自己。
  \item 选题兼顾兴趣、问题价值、资源与可行性。
  \item 写作重在把问题—方法—实验—贡献讲清楚。
  \item 尽早形成文献阅读、实验记录、阶段总结、定期汇报的习惯。
  \item 从真实问题出发，而不是为了发文而堆方法。
\end{enumerate}
